\section{Benchmark design and experimental protocol}
\label{sec:si-protocol}

\subsection{Task specification and study units}
CloudraBench comprises 50 tasks, distributed as 20 Easy, 10 Normal, 10 Hard, and 10 Impossible tasks. Easy tasks cover elementary transformations, conformational changes, and basic execution; Normal tasks require intermediate identification and continuity across elementary steps; Hard tasks emphasize competing pathways and stereochemical identity; Impossible tasks require conclusions appropriate to the capabilities of the allowed electronic-structure methods. A task fixes the reactant and target identities, reaction conditions, total charge, spin multiplicity, requested deliverables, and permitted methods. These specifications are common to the workflow configurations compared on that task. The complete task catalogue accompanies the per-task results.

The repeated-task study is identified as P1b and the full-benchmark study as P2. Two studies provide complementary experimental units (Table~\ref{tab:si-study-design}). The repeated-task study contains three independent runs for each combination of four tasks and four configurations. E1 concerns elimination preorganization; NB1 concerns multistep decarboxylation; HA1 concerns stereochemical ring opening; and IA1 concerns an electronic-structure method limitation. The full-benchmark study contains one retained run for each of the 50 tasks under each configuration. Its primary matched unit is therefore a task represented in all four configurations. The 48 repeated-task runs and the 200 full-benchmark runs are reported as separate studies.

\begin{table}[htbp]
\centering
\caption{Experimental units in the two studies.}
\label{tab:si-study-design}
\small
\begin{tabular}{p{0.21\textwidth}p{0.31\textwidth}p{0.34\textwidth}}
\toprule
Design & Repeated-task study (P1b) & Full-benchmark study (P2) \\
\midrule
Task set & E1, NB1, HA1, IA1 & All 50 CloudraBench tasks \\
Backbone & DeepSeek-v4.1-flash (1017) & DeepSeek-v4.1-flash (1017); official DeepSeek chat-completions API \\
Configurations & Guided, Standard, Record, Brief-errors & Guided, Standard, Record, Brief-errors \\
Repetitions & Three per task and configuration & One retained run per task and configuration \\
Task-level runs & $4\times4\times3=48$ & $50\times4=200$ \\
Elapsed-time unit & Four-task batch at concurrency two; three batches per configuration & Computational effort on the common 15-task subset \\
\bottomrule
\end{tabular}
\end{table}

\subsection{Workflow configurations and primary comparison}
Guided is compared primarily with the actual Standard branch using the same model backbone. Record and Brief-errors isolate changes to scientific acceptance and diagnostic presentation within the available workflow (Table~\ref{tab:si-configurations}). Standard retains access to computational tools while hiding explicit delivery-DAG tools, removing scientific acceptance interception, and providing neutral computational feedback. Guided retains the delivery graph, scientific checks, and detailed diagnostics. Record retains available tools and recording while removing mandatory scientific acceptance. Brief-errors retains guidance and acceptance while shortening diagnostics. The task and method restrictions are not relaxed in an ablation.

\begin{table}[htbp]
\centering
\caption{Operational definition of the four workflow configurations. Brief-errors is abbreviated as Brief in figures and tables; Standard is also labelled Backbone.}
\label{tab:si-configurations}
\small
\begin{tabular}{p{0.16\textwidth}p{0.25\textwidth}p{0.25\textwidth}p{0.22\textwidth}}
\toprule
Configuration & Delivery-graph interface & Mandatory scientific acceptance & Diagnostic presentation \\
\midrule
Guided & Explicit & Retained & Detailed \\
Standard & Explicit DAG tools hidden & Removed & Neutral computation feedback \\
Record & Available tools and records retained & Removed & Recording retained \\
Brief-errors & Retained & Retained & Shortened \\
\bottomrule
\end{tabular}
\end{table}

Both studies used DeepSeek-v4.1-flash (1017). The API alias is \texttt{deepseek-v4-flash} in all eight frozen coordinator and follower configurations and in their actual model-call telemetry.\citep{deepseekv42026} Model names, campaign identifiers, budgets, stopping reasons, and artifacts are retained with the corresponding runs. Cloudra revision 70daa5a identifies the implementation described in the main-text methods.\citep{ref14} This implementation identifier, the experiment records, and the scoring and organization records serve different purposes: software or API versions are taken from the applicable archive rather than inferred from a later installed environment or model alias.

\subsection{P2 model configuration and campaign versions}
\label{sec:si-p2-model-metadata}

P1b denotes the repeated-task study: four tasks, four workflow configurations, and three runs per combination, giving 48 task runs. P2 denotes the full-benchmark study of 50 tasks under the four configurations, giving 200 retained task runs. The P2b and P2c execution campaigns supply this P2 dataset. Three completed P2b runs were retained (Guided/D3, Record/D1, and Standard/D3), and P2c supplied the remaining 197 runs.

The deployed backbone was DeepSeek-v4.1-flash (1017), as specified by the authors. All eight P2 configuration snapshots, one for each campaign and workflow combination, specify \texttt{deepseek-v4-flash} for both the coordinator and follower slots. The request-start and completion telemetry in all eight runtime archives records the same model identifier. Cloudra called the official DeepSeek endpoint through streaming OpenAI-compatible chat completions. The frozen \texttt{native\_search=auto} setting selects this route in the campaign-era implementation. Table~\ref{tab:si-p2-model-metadata} records the request settings and execution versions. The deployment name and version are reported alongside the literal API alias preserved in the execution records.

\begin{table}[htbp]
\centering
\small
\caption{Frozen P2 configuration and campaign metadata. Request settings describe the harness configuration; telemetry counts are provided in the accompanying \texttt{p2\_model\_metadata.json}.}
\label{tab:si-p2-model-metadata}
\begin{tabular}{p{0.27\linewidth}p{0.66\linewidth}}
\toprule
Item & Recorded value \\
\midrule
Deployed model & DeepSeek-v4.1-flash (1017) \\
Provider and API alias & DeepSeek; \texttt{deepseek-v4-flash} in all eight snapshots and their request/completion telemetry \\
Endpoint and transport & \url{https://api.deepseek.com}; streaming chat completions \\
Role slots & Coordinator and follower use the same provider, model identifier, and request settings \\
Reasoning setting & Requested \texttt{reasoning\_effort=minimal} \\
Temperature & 1.0 in the archived run-start settings \\
Context window & 400,000 tokens in the frozen harness configuration \\
Spawned-agent budget & \texttt{max\_tokens=32768} in the archived \texttt{agent.spawned} settings; this is the per-agent request setting, not the run's total token consumption \\
Tool invocation & JavaScript protocol \\
Retained task runs & Three from P2b and 197 from P2c \\
Task concurrency & Two per configuration, eight in total \\
Campaign starts & P2b: 25 October 2026, 18:01:17 UTC+8; P2c: 25 October 2026, 22:18:45 UTC+8 \\
Campaign base & \texttt{53309d0e} \\
Configuration commits & Guided: \texttt{6e89d20c}; Record: \texttt{746c33f0}; Standard: \texttt{c7fc94a1}; Brief-errors: \texttt{3b9831bb}. Both campaigns use these commits. \\
\bottomrule
\end{tabular}
\end{table}

The configuration commits identify experiment execution. The implementation reference used to describe the system, subsequent evaluator fixes, and publication-table reproduction have their own version records. The companion metadata file gives archive-relative source members for each configuration and runtime record. Its request-start and completion counts refer to the archived runtime events, including auxiliary LLM activity and cancelled P2b attempts; the benchmark sample remains the 200 retained task runs.

\paragraph{Fixed geometric heuristics.}
The atom-map ranking uses fixed coefficients: 2 for each formed or broken bond, 1 for each bond-order change, 4 for a hydrogen-count mismatch, 8 for charge or radical-electron mismatches, and 2 for a stereochemical mismatch. The $J_{\min}+2$ window retains near-optimal maps when reporting ambiguity. The branch-crowding descriptor also uses fixed implementation heuristics: decay $\rho=0.5$, traversal through depth five, and a congestion flag when the four-branch total $B_{\mathrm{axis}}\geq9.0$. These settings rank or annotate candidate structures. Quantum-chemical evidence determines their subsequent physical qualification. The imaginary-mode participation ratio has a comparison value of 0.5; its alignment flag is advisory, and strict acceptance uses the first-order-saddle and IRC/endpoint evidence.

\paragraph{Conformational sampling.}
The default preorganization uses topology-aware rigid placements and guided MMFF94 relaxation, followed by energy/RMSD selection. Its standard multicomponent settings sample 32 relative orientations and retain up to eight guided placements. The separate conformer-search tool defaults to the RDKit backend and also provides an external CREST backend\citep{crest2020} for additional conformer or noncovalent-complex sampling. The default CREST route uses GFN2-xTB with the quick setting. Search outputs enter the calculation workflow as candidate geometries and undergo the task-specific stationary-point and pathway checks.



\subsection{Allowed calculations and reference definitions}
Easy, Normal, and Hard tasks permit xTB, PM7, and B3LYP-D3(BJ)/def2-SVP.\citep{gfn2xtb2019,pm72013,becke1993,lyp1988,stephens1994,d32010,d3bj2011,def22005} Impossible tasks are restricted to xTB and PM7. Native ORCA, Gaussian, and MOPAC records supply the corresponding computational quantities and termination information.\citep{orca2025,gaussian16,mopac2026,ref11} Each result is associated with the engine, calculation level, species, and task conditions represented in its source record.

The reference answer (Gold) specifies the mechanistic species and elementary steps, stereochemical identities, applicable numerical targets, and explicit negative targets. Numerical activation barriers use electronic energies at a consistent level of theory:
\begin{equation}
\Delta E_e^{\ddagger}=E_{\mathrm{TS},e}-E_{\mathrm{ref},e}.
\label{eq:si-barrier}
\end{equation}
These benchmark barriers do not include zero-point or thermal corrections. The energy reference is the verified stable minimum of the corresponding species. Declared conformational connections within the same species can share that stable reference; a conformational conversion itself retains its own starting conformer as the reference. Distinct stereochemical substrates and independent channel contact complexes retain separate references. The actual IRC connection and the energy denominator are recorded separately, so a stable reference is not automatically identified with the conformer reached directly by an IRC.

Reference treatment follows the individual task definition. NB1\_E2 and NQ5\_E2 have no numerical reference under the specified proton-transfer policy. IB9 combines a numerical first-step reference of 5.65~kcal~mol$^{-1}$ with a negative second-step target. The qualitative selectivity tasks and numerical tolerance settings are specified in Section~\ref{sec:si-scoring}. Reference status belongs to the corresponding reference record; task availability, presence of a Gold graph, and physical acceptance of a complete reference chain are recorded separately.

\subsection{Physical evidence associated with a mechanism}
A delivered mechanism represents species occurrences as nodes and elementary steps as edges. A species minimum requires convergence and a frequency calculation with zero imaginary frequencies. An accessible transition-state step requires a first-order saddle with exactly one imaginary frequency associated with the intended reaction, together with direction evidence linking that saddle to the declared endpoint species.\citep{ref09,ref10} Native normal termination, geometry convergence, frequencies, and pathway connectivity are checked as separate facts.

The direction evidence starts at the same transition state in the two IRC directions. The actual outer frame from each direction is connected through a converged optimization to an endpoint minimum with zero imaginary frequencies and the appropriate chemical identity and stereochemistry. The IRC trajectory itself need not converge all the way to either minimum when this outer-frame--optimization--frequency chain establishes the connection. The archived frames and their calculation lineage are retained. Geometric interval measurements assist interpretation; they do not replace a supported native endpoint chain. For separated products, the required constituent minima are verified; an intact molecular complex is evaluated as the corresponding whole-system species.

These requirements apply to the objects requested by the task. An explicitly negative target is represented as inaccessible with its basis retained. A species-only negative question does not require an invented transition state, process edge, or numerical barrier.

\subsection{Offline organization and publication records}
Formal evaluation is conducted after the run artifacts are organized into the common mechanism-graph representation. Organization uses that run's conversations, tool records, reports, geometries, and native calculation outputs to recover graph objects, repair file references, and convert quantities already present in the run. Human organization and field-level sources are recorded, and the same organizational rules apply across configurations. The resulting graph expresses the scientific content of the run; Gold is used for evaluation rather than as a source of additional run conclusions or computational evidence.

The full-benchmark publication table contains 200 runs and all 600 Graph, Self-consistency, and Accuracy component values. Each published composite is aligned with its own existing component records. The 200 released composite scores are retained; the weighted components reproduce them to their four-decimal reporting precision. The publication tables separately preserve whether the run itself generated a root mechanism graph and whether the graph was subsequently organized offline. Thus, the original delivery behavior and the organized scientific score remain separately inspectable quantities.

The repeated-task study retains its historical four-configuration snapshots. The subsequent Standard offline organization is a separately identified update to those records; it is not labelled as a new uniform rescore of all 48 runs. Analyses using the repeated-task archive specify which record supplies the result.

\subsection{Paired statistical analysis and time measurements}
For a task $i$ and configuration $a$, let $S_{ia}$ denote its composite score. Configuration means give each task equal weight within the reported set. The primary paired difference is
\begin{equation}
\overline{\Delta}_{\mathrm{G-S}}=\frac{1}{n}\sum_{i=1}^{n}
\left(S_{i,\mathrm{Guided}}-S_{i,\mathrm{Standard}}\right).
\label{eq:si-paired-difference}
\end{equation}
The same calculation is used within each tier and for the combined Normal/Hard subset. Weighted component contributions are the mean paired component differences multiplied by 0.1, 0.3, and 0.6, respectively. They add to the composite difference up to the reporting precision of the released scores.

The 95\% percentile intervals use 100,000 task-cluster bootstrap samples with seed 20261031. Each resampled task carries its four configuration records together, preserving the paired comparison. The bootstrap unit is a task, not an elementary edge or an individual quantum-chemical calculation.

The repeated-task elapsed-time analysis uses the recorded start and finish of each four-task batch at concurrency two: three batches per configuration, or 12 batches in total. Its unit is batch wall-clock hours. The full-benchmark effort analysis uses the 15 tasks with common effort records across all four configurations, comprising 60 task-level runs. It reports calculation counts and cumulative quantum-chemical job-hours, obtained by summing the individual calculation durations within a run. In the archived Standard/M1 extraction, native-log CPU time is divided by four to obtain the corresponding job-duration estimate. Job-hours are neither core-hours nor a measurement of the full P2 research-run wall-clock duration. These time quantities retain their distinct denominators and units throughout the supplementary tables.
